\begin{description}
\item[(2.1)] Geographical latitude of the observing site: $25^\circ$,
  Southern hemisphere. The constellation Crux shows the direction of the
  southern celestial pole; the altitude of the pole gives the latitude.
  (Latitude between $24^\circ$--$26^\circ$ = 2 points, within $\pm 3^\circ$
  = 1 point; Southern hemisphere = 1 point.) \hfill(maximum 3 points)
\item[(2.2)] % note: the official solution measures azimuth from South
  % towards West (0--360), unlike the answer-sheet wording
  1st brightest star: Sirius ($\alpha$ CMa), $A = 82^\circ$
  ($74^\circ$--$90^\circ$ = 2 points, $64^\circ$--$74^\circ$ and
  $90^\circ$--$100^\circ$ = 1 point);\\
  2nd brightest star: Canopus ($\alpha$ Car), $A = 42^\circ$
  ($34^\circ$--$50^\circ$ = 2 points, $24^\circ$--$34^\circ$ and
  $50^\circ$--$66^\circ$ = 1 point);\\
  3rd brightest star: Toliman ($\alpha$ Cen), $A = 331^\circ$
  ($324^\circ$--$339^\circ$ = 2 points, $314^\circ$--$324^\circ$ and
  $339^\circ$--$349^\circ$ = 1 point).\\
  (Star name 2 points each; incorrect position in the list $-1$ point.)
  \hfill(maximum 12 points)
\item[(2.3)] The number two comet is in the closest position to the
  ecliptic. (Distances of the comets from the ecliptic: No.\ 1 =
  $20^\circ$, No.\ 2 = $0^\circ$, No.\ 3 = $30^\circ$.) The line
  connecting Regulus and Spica approximately shows the direction of the
  ecliptic. \hfill(2 points)
\item[(2.4)] 1 point for each of: Octans (Oct), Mensa (Men), Chamaeleon
  (Cha), Apus (Aps), Hydrus (Hyi), Volans (Vol), Carina (Car), Musca
  (Mus), Circinus (Cir), Triangulum Australe (TrA), Pavo (Pav), Indus
  (Ind), Tucana (Tuc), Horologium (Hor), Ara (Ara).
  \hfill(maximum 9 points)
\item[(2.5)] 12 hours. $\delta$ Orionis is located approximately on the
  celestial equator ($\delta = -0^\circ 18'$). \hfill(2 points)
\end{description}
